%%%%%%%%%%%%%%%%%%%%%%%%%%%%%%%%%%%%%%%%%%%%%%%%%%%%%
%			HLAVIČKA								%
%%%%%%%%%%%%%%%%%%%%%%%%%%%%%%%%%%%%%%%%%%%%%%%%%%%%%
\documentclass[openany,12pt]{memoir}

\usepackage[utf8]{inputenc} 
\usepackage[czech]{babel}
\usepackage[T1]{fontenc}
\usepackage[top=1.5cm, bottom=2cm, left=2cm, right=2cm]{geometry}  % --> NASTAVENÍ OKRAJŮ
\usepackage{fancyhdr}
\usepackage{graphicx}
%\usepackage{lmodern}
\usepackage{tgschola}  % --- Silnější font
\usepackage{xwatermark}
\usepackage{xcolor}
\usepackage{changepage}
\usepackage{newunicodechar}



%%%%%% Package na zpěvník
\usepackage[full]{leadsheets}%http://mirrors.nic.cz/tex-archive/macros/latex/contrib/leadsheets/leadsheets_en.pdf   --> dokumentace	
\definesongtitletemplate{empty}{} 
\setchords{
format = \bfseries,   %tučné akordy
minor = {mi},% 
input-notation = {german},%
output-notation = {german}%
}
\definesongtitletemplate{empty}{} 

\newlength{\drop}
\newwatermark[allpages,color=red!50,angle=0,scale=2, xpos=0,ypos=0]{\includegraphics[width=5cm]{obr/pozadi2.jpg}} %--> dvojka na pozadí


%%%% Vlastní příkazy
\newcounter{Slokočet}   %Automatické číslování slok
\newcommand{\mezera}{\vspace*{0.5cm}}   %Horizontální odsazení slok
\newcommand{\stred}{5.2cm}   %%% Na zarovnání slok doprostřed, pozn. automatičtější zarovnávání na střed nejde
\newcommand{\refren}{\mezera \noindent \textbf{R:} } %refrén
\newcommand{\sloka}{\mezera \noindent \addtocounter{Slokočet}{1} \arabic{Slokočet}. } 	%sloka, která se automaticky čísluje
\newcommand{\ssloka}{\mezera \noindent}  % vlastní číslo sloky
\newcommand{\carka}{,\:}
\newcommand{\predehra}{\mezera \hspace*{-1.85cm} \textbf{Předehra:} }
\newcommand{\m}[1]{\color{white}{#1}}  %Pro akordy
\newcommand{\ap}{'} %Pro apostrof
\newcommand{\elipsa}{\kern\fontdimen3\font} %Příkaz pro lepší zacházení s výpustkami (=...); je to vpodstatě jen mezera mezi tečkama výpustky

%%%%%%%%%%%%%%%%%%%%%%%%%%%%%%%%%%%%%%%%%%%%%%%%%%%%%%
%			NÁVOD									 %
%%%%%%%%%%%%%%%%%%%%%%%%%%%%%%%%%%%%%%%%%%%%%%%%%%%%%%
%1. Věci v hlavičce IGNOROVAT
%2. Píseň psát do prostoru mezi \begin{song} a \end{song}
%3. další řádek se značí dvěma odsazeníma (= dvakrát stisknout enter)
%4. \refren vždy na začátku refrenu a \sloka na začátek sloky (automaticky se čísluje)
%  \ = alt gr + q ; [/] = alt gr f/g ; {/} = alt gr + b/n; ^ = alt gr + 3 + mezera
%Cokoliv napíšete do ^{  } se bude brát jako akord
%když se toto bude dotýkat nějakého slova (nebude mezi tím a slovem mezera)
%tak se akordy zjeví nad slovem, ale pište to před slovo
%Když se to nedotýká slova, tak akord lítá ve vzduchu a vytiskne se větší mezera
%První možnost je asi preferovanější
%5. Akordy stačí psát jen do první sloky, když se nezmění -- kytaristi to zvládnou



%\usepackage{subfiles}
%</++>

\begin{document}

\pagestyle{empty}
%\addtocontents{toc}{\protect\thispagestyle{empty}}
%\tableofcontents \thispagestyle{empty}\newpage
\newgeometry{top=1.5cm, bottom = 0cm, left = 2cm, right = 2cm}

\newpage

%Pro kompilaci je důležitý, 	aby ten input a \end{document} zůstali na posledních
%dvou řádcích!!!!
\begin{song}{title=\centering Nebeští Jezdci \\\normalsize \vspace*{-0.3cm}}  %% sem se napíše jméno songu a autor
\moveright -0.9cm \vbox{      %Varianta č. 1  ---> Jeden sloupec zarovnaný na střed	
\begin{minipage}[t]{0.57\textwidth}\setlength{\parindent}{0.45cm}  %Varianta č. 2 --> Dva sloupce

\sloka
	Jel ^{Emi{\color{white}\_}}starý cowboy, jel ^{G{\color{white}\_\_}}stále jel a jel,
	
	jen ^{Emi{\color{white}\_}}jeden cíl měl před sebou, ten ^{G{\color{white}\_\_}}dávno ^{{\color{white}\_\_}Emi}zapomněl,
	
	^{C{\color{white}\_}}když tu náhle ^{Emi\,\,}jeho kůň u ^{C{\color{white}\_\_}}skály zůstal ^{Emi}stát,
	
	^{C}co starý cowboy uviděl, to ^{Emi{\color{white}\_}}nelze vyzpívat.
	
	Yupy ya^{G}hej, yupa ya ^{Emi}hou,
	
	^{C}co starý cowboy uviděl, to ^{Emi{\color{white}\_}}nelze vyzpívat.
	
\sloka
	Kde rudé slunko zapadá, tam cowboy jezdce zří,
	
	a neskonalým štěstím se mu oči rozzáří,
	
	jednou na bílých koních, jak je to nádherné,
	
	podkovy mají zlaté a uzdy stříbrné.
	
	Yupy ya hej, yupa ya hou,
	
	podkovy mají zlaté a uzdy stříbrné.
	
\sloka
	Tam vpředu jede muž, který již dávno mrtvý byl,
	
	a za ním jede ten, co za něj život položil.
	
	\uv{Počkejte na mne kamarádi chci jet s Vámi dál!}
	
	Když ho našli mrtvého, ještě se usmíval.
	
	Yupy ya hej, yupa ya hou,
	
	když ho našli mrtvého, ještě se usmíval.
	
\sloka
	Nebeští jezdci čekejte, chci ženě sbohem dát,
	
	že s Vámi jsem, že nemusí se o mne vůbec bát,
	
	že vrátím se na bílém koni zase jednou k ní,
	
	bude to má dlouhá, má jízda poslední.
	
	Yupy ya hej, yupa ya hou,
	
	bude to má dlouhá, má jízda poslední.
		

\end{minipage}\begin{minipage}[t]{0.55\textwidth}\setlength{\parindent}{0.45cm}\vspace*{0.55cm}  % V případě varianty č.2 jde odsud text do pravé části

\sloka
	Nebeští jezdci slyším Vás a chce se mi tak spát,
	
	tu velký muž na bílém koni v sedle zůstal stát.
	
	\uv{Ty nevrátíš se více, ať to každý ví,
	
	bude to tvá jízda, ta jizda poslední.}
	
	Yupy ya hej, yupa ya hou,
	
	bude to tvá jízda, ta jizda poslední.
	
\sloka
	Nebeští jezdci jedou dál a vítr přestal vát,
	
	ten velký muž mi pokynul, já na cestu se dal,
	
	já nevrátím se více, ať to každý ví,
	
	je to moje dlouhá, má jízda poslední.
	
	Yupy ya hej, yupa ya hou,
	
	je to moje dlouhá, má jízda poslední.


\end{minipage}
}
\setcounter{Slokočet}{0}
\end{song}


\end{document}
